\documentclass[12pt]{article}
\usepackage[margin=1in]{geometry}
\usepackage{amsmath}
\usepackage{amssymb}
\usepackage{amsthm}
\usepackage{enumitem}
\usepackage{xcolor}
\usepackage{pagecolor}
\usepackage{marginnote}
\usepackage{fancyhdr}
\usepackage{bm}
\usepackage{etoolbox}
\usepackage{mdframed}

\definecolor{cream}{HTML}{faf7f1}
\pagecolor{cream}
\mdfdefinestyle{lemmastyle}{backgroundcolor=cream}

\newcommand{\defn}[1]{\textbf{#1}}
\newcommand{\defeq}{\mathrel{\vcenter{\baselineskip0.5ex \lineskiplimit0pt
                     \hbox{\scriptsize.}\hbox{\scriptsize.}}}=}
\newcommand{\T}{\top}
\newcommand{\F}{\bot}
\newcommand{\Pow}{\mathbb{P}}
\newcommand{\suc}{\mathfrak{suc}}
\newcommand{\N}{\mathbb{N}}
\newcommand{\Z}{\mathbb{Z}}
\newcommand{\Q}{\mathbb{Q}}
\newcommand{\Func}{\mathcal{F}}
\newcommand{\id}{\mathrm{id}}

\newtheoremstyle{proofstyle}
  {}{}{}{}{\bfseries}{.}{ }{}
\theoremstyle{proofstyle}
\newtheorem*{proof*}{Proof}

\newcommand{\sidenote}[1]{\marginnote{\small\itshape #1}}

\makeatletter
\renewcommand{\maketitle}{%
  \begin{flushleft}
    {\Large\bfseries CS173: Discrete Structures}\\[0.3em]
    {\Large\bfseries Problem Set 8}\\[0.5em]
    {\normalsize Mohammed Al Salhi}
  \end{flushleft}
  \vspace{1em}
}
\makeatother

\AtBeginDocument{\mathversion{bold}}

\begin{document}

\maketitle

\begin{enumerate}[left=0pt, itemsep=2em, parsep=1em]

%% ---------------------------------------------------------------
%% Problem 1
%% ---------------------------------------------------------------
\item An archaeologist on a recent expedition to the British Isles has discovered, for each
positive natural number $n \in \mathbb{N}^{+}$, evidence of a manuscript written in a language
called nglish. Each such manuscript is exactly $n$ pages long and contains precisely $2n$
distinct words. Prove each manuscript contains a page with two distinct words on it.

\begin{proof*}

Let $n \in \mathbb{N}^{+}$ be arbitrary, and consider an nglish manuscript with $n$ pages
and $2n$ distinct words. For each $i \in n$, let $P_i$ denote the set of words on page $i$.

Since the manuscript contains $2n$ distinct words,
$\left|\bigcup_{i=0}^{n-1} P_i\right| = 2n$.

Assume for contradiction that no page contains two distinct words, i.e.,
$(\forall i \in n)(|P_i| \leq 1)$. Observe:
\[
  2n = \left|\bigcup_{i=0}^{n-1} P_i\right| \leq \sum_{i=0}^{n-1} |P_i| \leq \sum_{i=0}^{n-1} 1 = n.
\]
This gives $2n \leq n$, a contradiction since $n \in \mathbb{N}^{+}$.

Therefore $(\exists i \in n)(|P_i| \geq 2)$, i.e., a manuscript contains a page with two
distinct words. \qed

\end{proof*}

\newpage

%% ---------------------------------------------------------------
%% Problem 2
%% ---------------------------------------------------------------
\item Show that there exist $a, b \in \mathbb{N}$ such that $256 < a < b$ and $7^{a} \equiv 7^{b} \pmod{512}$.

\begin{proof*}
By the Euclidean Division Lemma,
\[
(\forall a \in \mathbb{Z})(\forall b \in \mathbb{Z} \setminus \{0\})(\exists! q, r \in \mathbb{Z})(a = bq + r \land 0 \leq r < |b|).
\]
Define $\mathbin{\%} : \mathbb{Z} \times (\mathbb{Z} \setminus \{0\}) \to \mathbb{Z}$ by
$x \mathbin{\%} y \defeq r$, where $r$ is the unique integer guaranteed by the Euclidean
Division Lemma satisfying $x = yq + r$ and $0 \leq r < |y|$ for some $q \in \mathbb{Z}$.

Let $A \defeq \mathbb{N} \setminus \{k \in \mathbb{N} \mid k < 257\}$ and
$B \defeq \{k \in \mathbb{N} \mid k < 512\}$, so $|B| = 512$.

Define $f : A \to B$ by $(\forall x \in A)(f(x) \defeq 7^x \mathbin{\%} 512)$.

By the Corollary, $|A| = |\mathbb{N}|$. By the Theorem, an infinite set is larger than any
finite set, so
\[
|A| = |\mathbb{N}| > |B| = 512.
\]

By the Pigeonhole Theorem,
\[
    |A| > |B| \Rightarrow (\forall f : A \to B)(\exists x, y \in A)(x \neq y \land f(x) = f(y)),
\]
so $(\exists c, d \in A)(c \neq d \land f(c) = f(d))$, i.e., $7^c \mathbin{\%} 512 = 7^d \mathbin{\%} 512$.

By the Euclidean Division Lemma, $7^c = 512 q_1 + r_1$ and $7^d = 512 q_2 + r_2$ with
$0 \leq r_1, r_2 < 512$. Since $f(c) = f(d)$, we have $r_1 = r_2$, so subtracting gives
$7^c - 7^d = 512(q_1 - q_2)$. Since $q_1 - q_2 \in \mathbb{Z}$, we have
$512 \mid (7^c - 7^d)$, i.e., $7^c \equiv 7^d \pmod{512}$.

Let $a \defeq \min(c, d)$ and $b \defeq \max(c, d)$. Since $c \neq d$ we have $a < b$,
and since $c, d \in A$ we have $a, b > 256$.

Therefore there exist $a, b \in \mathbb{N}$ such that $256 < a < b$ and $7^a \equiv 7^b \pmod{512}$. \qed
\end{proof*}

\newpage

%% ---------------------------------------------------------------
%% Problem 3
%% ---------------------------------------------------------------
\item Consider the set $\mathcal{S} \defeq \{3, 4, 7, 8, 9, 10, 12, 15, 18, 19, 27, 28\}$. Show that, in any subset
$\mathcal{X} \subseteq \mathcal{S}$ of cardinality $|\mathcal{X}| \geq 9$, there exist three distinct elements
$x_1, x_2, x_3 \in \mathcal{X}$ such that $x_1 + x_2 + x_3 = 40$.

\begin{proof*}
Partition $\mathcal{S}$ into four triples:
\[
  G_1 = \{3, 10, 27\}, \quad G_2 = \{9, 12, 19\}, \quad G_3 = \{4, 8, 28\}, \quad G_4 = \{7, 15, 18\},
\]
where $3+10+27 = 9+12+19 = 4+8+28 = 7+15+18 = 40$, the $G_i$ are pairwise disjoint,
and $G_1 \cup G_2 \cup G_3 \cup G_4 = \mathcal{S}$.

Let $\mathcal{X} \subseteq \mathcal{S}$ with $|\mathcal{X}| \geq 9$. Suppose for contradiction
that no three distinct elements of $\mathcal{X}$ sum to 40. If $G_i \subseteq \mathcal{X}$
for some $i$, then the three distinct elements of $G_i$ are in $\mathcal{X}$ and sum to 40,
contradicting our assumption. Hence $|\mathcal{X} \cap G_i| \leq 2$ for each $i$.

Since $\{G_1, G_2, G_3, G_4\}$ partitions $\mathcal{S}$ and $\mathcal{X} \subseteq \mathcal{S}$,
the sets $\mathcal{X} \cap G_1, \ldots, \mathcal{X} \cap G_4$ are pairwise disjoint with union
$\mathcal{X}$. Observe,
\[
  |\mathcal{X}| = \sum_{i=1}^{4} |\mathcal{X} \cap G_i| \leq 4 \times 2 = 8,
\]
contradicting $|\mathcal{X} \geq 9|$.

Therefore our claim that $|\mathcal{X} \cap G_i| \leq 2$ for all $i$ is false, so there
exists some $i$ with $|\mathcal{X} \cap G_i| = 3$, i.e., $G_i \subseteq \mathcal{X}$. The
three distinct elements of $G_i$ are then in $\mathcal{X}$ and sum to 40, giving
$x_1, x_2, x_3 \in \mathcal{X}$ with $x_1 + x_2 + x_3 = 40$. \qed
\end{proof*}
\newpage

%% ---------------------------------------------------------------
%% Problem 4
%% ---------------------------------------------------------------
\item For this problem, you may assume that $\binom{n}{0} = \binom{n}{n} = 1$ for all $n \in \mathbb{N}$.
Prove that, for every $n, k \in \mathbb{N}$ such that $k \leq n$, we have
$\binom{n}{k} = \binom{n}{n-k}$.

\begin{proof*}
Let $n, k \in \mathbb{N}$ with $k \leq n$. Define
\[
  S_k \defeq \{B \mid B \subseteq n \land |B| = k\}, \qquad
  S_{n-k} \defeq \{B \mid B \subseteq n \land |B| = n-k\},
\]
so that $\binom{n}{k} = |S_k|$ and $\binom{n}{n-k} = |S_{n-k}|$ by definition.

Define $f : S_k \to S_{n-k}$ by $(\forall B \in S_k)(f(B) \defeq n \setminus B)$.

\textbf{Well-defined.} Let $B \in S_k$. Then $f(B) = n \setminus B \subseteq n$, and by the
Theorem ($B \subseteq n \Rightarrow |n \setminus B| = |n| - |B|$), we have
$|n \setminus B| = n - k$, so $f(B) \in S_{n-k}$.

\textbf{Injective.} Let $B_1, B_2 \in S_k$ with $f(B_1) = f(B_2)$. Then
$n \setminus B_1 = n \setminus B_2$. Since $B_1, B_2 \subseteq n$, we have
$B_1 = n \setminus (n \setminus B_1) = n \setminus (n \setminus B_2) = B_2$.

\textbf{Surjective.} Let $C \in S_{n-k}$. Then $n \setminus C \subseteq n$, and by the
Theorem, $|n \setminus C| = |n| - |C| = n - (n-k) = k$, so $n \setminus C \in S_k$.
Then $f(n \setminus C) = n \setminus (n \setminus C) = C$.

Since $f$ is a bijection, $|S_k| = |S_{n-k}|$, i.e., $\binom{n}{k} = \binom{n}{n-k}$. \qed
\end{proof*}
\newpage

%% ---------------------------------------------------------------
%% Problem 5
%% ---------------------------------------------------------------
\item We define a \defn{finite binary string} to be a function $b : n \to \{0, 1\}$, where $n \in \mathbb{N}$
is the length of the string. We define an \defn{infinite binary string} to be a function
$b : \mathbb{N} \to \{0, 1\}$.

How many infinite binary strings are there that each contain one contiguous finite
substring consisting only of $1$s? Formalize your claim and justify it with a proof.

\textbf{Claim:} There are countably many infinite binary strings containing exactly one
contiguous finite substring of all 1s, i.e., the set of such strings has cardinality $|\mathbb{N}|$.

\begin{proof*}
Let $\mathcal{T}$ denote the set of all infinite binary strings $b : \mathbb{N} \to \{0,1\}$
such that the set $\{n \in \mathbb{N} \mid b(n) = 1\}$ is a nonempty finite contiguous
interval.

For $i, k \in \mathbb{N}$ with $i \leq k$, define $b_{i,k} : \mathbb{N} \to \{0,1\}$ by
\[
  b_{i,k}(n) \defeq \begin{cases} 0 & \text{if } n < i \\ 1 & \text{if } i \leq n \leq k \\ 0 & \text{if } n > k \end{cases}
\]

Let $S \defeq \{(i,k) \in \mathbb{N} \times \mathbb{N} \mid i \leq k\}$ and define
$f : S \to \mathcal{T}$ by $f(i,k) \defeq b_{i,k}$.

\textbf{Injective.} Let $(i_1, k_1), (i_2, k_2) \in S$ with $f(i_1, k_1) = f(i_2, k_2)$,
and let $b = b_{i_1, k_1} = b_{i_2, k_2}$. Then
$i_1 = \min\{n \in \mathbb{N} \mid b(n) = 1\} = i_2$ and
$k_1 = \max\{n \in \mathbb{N} \mid b(n) = 1\} = k_2$.

\textbf{Surjective.} Let $b \in \mathcal{T}$. Since $b$ has exactly one maximal block of 1s,
let $i \defeq \min\{n \in \mathbb{N} \mid b(n) = 1\}$ and $k \defeq \max\{n \in \mathbb{N} \mid b(n) = 1\}$.
Then $(i, k) \in S$ and $f(i, k) = b$.

Since $f$ is a bijection, $|S| = |\mathcal{T}|$. The set $\{(0, k) \mid k \in \mathbb{N}\} \subseteq S$
is infinite, so $S$ is infinite. Since $S \subseteq \mathbb{N} \times \mathbb{N}$ and
$|\mathbb{N} \times \mathbb{N}| = |\mathbb{N}|$, by Theorem, so $|S| = |\mathbb{N}|$. Therefore $|\mathcal{T}| = |\mathbb{N}|$.

Including the all-zero string adds one element to our count, giving $|\mathcal{T}| + 1$
total strings. Since $\mathcal{T}$ is infinite, adding a single element does not change
its cardinality, so the total count remains $|\mathbb{N}|$. \qed
\end{proof*}
\end{enumerate}

\end{document}